% ============================================================================
% preamble/06-thesis.tex -- what this thesis needs on top of the template
%
% preamble/01..05 are the institute template's own preamble files and are kept
% as close to upstream as the switch to English allows, so that they stay
% diffable against ~/latex-template-thesis. Everything this thesis adds lives
% here.
% ============================================================================

% ----- Figures --------------------------------------------------------------
% Figures are \input from the chapter that discusses them and reference images
% by bare filename.
\graphicspath{{figures/}}

% The template sets no caption fonts, so KOMA's defaults apply. subcaption
% loads caption, which would otherwise take over; tell it to keep using them.
\captionsetup{font=small, labelfont=bf}
\captionsetup[sub]{font=small, labelfont=bf}

% ----- Diagrams -------------------------------------------------------------
% preamble/03-graphics.tex already loads tikz with positioning, calc, fit and
% friends; the automata drawings need these two on top.
\usetikzlibrary{arrows.meta,automata}
\tikzset{terminal/.style={double,double distance=1.2pt}}

% ----- Result plots ---------------------------------------------------------
% The figures under figures/generated/ are written by scripts/evaluate.py out of
% the campaign databases and read their numbers from figures/data/*.csv. They
% need two pgfplots libraries the template does not load: fillbetween, for the
% quantile bands every distribution is drawn as, and colormaps for the named
% maps the collapse heatmaps use.
\usepgfplotslibrary{fillbetween}
\usepgfplotslibrary{colormaps}

% The generated figures set only what distinguishes them from one another;
% what they share lives here, so restyling the results does not mean
% regenerating them. Note the axis lines: the template's default is `center',
% which runs the axes through the origin. That is right for a function plot and
% wrong for a measurement read against a non-zero baseline, which is what every
% result plot in this thesis is.
% The label anchors need setting explicitly. `axis lines=center' installs label
% styles that park the labels at the far ends of the axes, which is where they
% belong when the axes cross at the origin; because the template sets that
% globally, the styles survive a per-axis switch to `axis lines=left' and the
% labels end up hanging off the corners of the plot. Re-centring them on the
% tick label band is what `axis lines=left' would have done on its own.
%
% The `rotate=90' on the ordinate label is load-bearing, not decoration: an
% unrotated label is as wide as its text, pgfplots does not count that width
% against the axis `width', and three panels of it overflow the text block.
% `scale only axis' (set by the template) makes `width' the width of the plot
% area alone, so labels, tick labels and colour bars stick out beyond it and a
% picture sized to its container overflows. Turning it off for these figures
% makes `width' the width of the whole picture, which is what a generated
% figure sized in \linewidth needs in order to fit its column.
\pgfplotsset{
  thesis result plot/.style={%
    axis lines=left,%
    axis on top=false,%
    scale only axis=false,%
    tick align=outside,%
    grid=major,%
    grid style={KITblack15,very thin},%
    legend cell align=left,%
    xlabel style={at={(ticklabel cs:0.5)},anchor=near ticklabel},%
    ylabel style={at={(ticklabel cs:0.5)},rotate=90,anchor=near ticklabel}},%
  thesis matrix plot/.style={%
    axis lines=left,%
    scale only axis=false,%
    tick align=outside,%
    enlargelimits=false,%
    axis on top,%
    xlabel style={at={(ticklabel cs:0.5)},anchor=near ticklabel},%
    ylabel style={at={(ticklabel cs:0.5)},rotate=90,anchor=near ticklabel}},%
}

% ----- Tables ---------------------------------------------------------------
% tabularx also arrives via the template's ltxtable; loading it by name states
% the dependency rather than relying on that.
\usepackage{tabularx}

% ----- Lists ----------------------------------------------------------------
% Provided by the previous template (sdqthesis.cls). The research questions in
% sections/introduction.tex use its label/ref keys, which plain enumerate has
% no equivalent for.
\usepackage{enumitem}

% ----- Units & numbers ------------------------------------------------------
\usepackage{siunitx}

% ----- Conditionals ---------------------------------------------------------
% \ifthenelse, used by the title page (an advisor line that is only printed
% when there is a second advisor) and by chapter.tex.
\usepackage{ifthen}

% ----- Algorithms -----------------------------------------------------------
% float, which algorithm needs for [H], is already loaded by the template's
% rotfloat.
\usepackage{algorithm}
\usepackage{algpseudocode}

% ----- Text macros ----------------------------------------------------------
% Provided by the previous template (sdqthesis.cls) and used throughout the
% text: \code for identifiers and paths, \model for model element names.
\newcommand{\code}[1]{\texttt{\hyphenchar\font45\relax #1}}
\newcommand{\model}[1]{\textsf{#1}}

% ----- Network names --------------------------------------------------------
% The EfficientZero networks follow MuZero's notation --- h (representation),
% g (dynamics), p (policy), v (value), vp (value prefix) --- with the parameter
% vector as an explicit argument, v(s;\theta), rather than a subscript. The
% multi-letter names need \mathit: plain math mode would set them as a product
% of single-letter variables, with the wrong spacing.
\newcommand{\vpnet}{\mathit{vp}}
\newcommand{\projnet}{\mathit{proj}}
\newcommand{\prednet}{\mathit{pred}}

% ----- Draft aids -----------------------------------------------------------
% Follows the draft switch set in preamble.tex: notes are shown while writing
% and hidden in the final document. The margin this template leaves is a little
% under todonotes' 2cm minimum. Widening \marginparwidth changes nothing in the
% final document --- \marginpar is only ever used by todonotes, which is
% disabled there.
\setlength{\marginparwidth}{2cm}
\ifthesisdraft
  \usepackage[colorinlistoftodos]{todonotes}
\else
  \usepackage[disable]{todonotes}
\fi
